Define the global forcing floor
\begin{equation}\label{eq:global-floor}
 \mathcal F_0=
 \min\{R_w(c_2,k_0),\Pi_{k_0}\}.
\end{equation}
After increasing $k_0$, monotonicity gives
$R_w(c_2,k_0)\le R_w(c_2,k)$ and
$R_w(c_2,k_0)\le\Pi_k$ for every relevant block and boundary.

\begin{lemma}[Localization]\label{lem:localization}
For sufficiently large $k_0$, every $a\notin\mathfrak M(C)$ satisfies either
\[
 \Qctrl(a)\ge\mathcal F_0
\]
or there is one integer $m$ with
\begin{equation}\label{eq:diagonal-sector}
 a_p\equiv m\pmod p\quad(p\in\Bset),
 \qquad |m|>C/\sctrl,
 \qquad \Qctrl(a)=m^2\sctrl^2.
\end{equation}
\end{lemma}
\begin{proof}
Suppose $\Qctrl(a)<\mathcal F_0$.  Then no block is hot and no adjacent cold
labels differ.  The small-$c_2$ refinement following
\eqref{eq:exception-count} eliminates every exception, so all coordinates
have one common label $m$.  The final scale also gives $2|m|<pq$ for every
control pair.  Hence $m$ itself is the centered CRT representative and
\[
 \Qctrl(a)=\sum_{(p,q)\in\Cgraph}(m/pq)^2=m^2\sctrl^2.
\]
If $|m|\le C/\sctrl$, then $a\in\mathfrak M(C)$; otherwise
\eqref{eq:diagonal-sector} holds.
\end{proof}

